\documentclass[10pt]{beamer}
\usepackage[utf8]{inputenc}
\usepackage[T1]{fontenc}
\usepackage{lmodern}
\usepackage{graphicx}
\usepackage{booktabs}
\usepackage{bm}
\usepackage{tikz}
\usepackage{ulem}
\usepackage{appendixnumberbeamer}
\usepackage{hyperref}
\usepackage{ragged2e}
\setbeamersize{text margin left=0.5cm, text margin right=0.5cm}
\definecolor{unigreen}{RGB}{3,138,94}
\usecolortheme[named=unigreen]{structure}
\beamertemplatenavigationsymbolsempty

\title{AI Workflows for Empirical Research}
\subtitle{General Setup: Claude Code \& Codex}
\author{Timur Öztürk}
\institute{University of Bayreuth}
\date{\today}

\begin{document}

\begin{frame}
  \titlepage
\end{frame}

\begin{frame}{The instructions file}
  \begin{itemize}
    \item Agentic tools read a Markdown file at every session start, and with every prompt. Keep it simple.
    \item Codex calls it \texttt{AGENTS.md}; Claude Code calls
          it \texttt{CLAUDE.md} (and also reads \texttt{AGENTS.md}).
    \item Just Markdown, no required fields. It can practically be anything.
    \item Like a ``README for the agent'': what you'd tell a new RA, written once.
  \end{itemize}
  \vfill
  \href{../sample_agents.md/}{\beamergotobutton{agents.md}}
  \href{https://code.claude.com/docs/en/memory}{\beamergotobutton{Claude Code memory}}
\end{frame}

\begin{frame}{Using it well}
  \begin{itemize}
    \item Keep it short (Anthropic: $<\sim$200 lines); put run/test commands at the top.
    \item Sections: overview, setup commands, conventions (e.g.\ cluster SEs),
          data-security notes.
    \item \textbf{MCPs} = connections to external tools/data;
          \textbf{skills} = reusable, on-demand instructions.
    \item Both run in a sandbox and can execute code!
  \end{itemize}
  \vfill
  \href{https://codersera.com/blog/agents-md-complete-guide-2026/}{\beamergotobutton{AGENTS.md guide}}\quad
  \href{https://dev.to/nishilbhave/claudemd-best-practices-the-complete-2026-guide-435j}{\beamergotobutton{CLAUDE.md best practices}}\quad
  \href{https://www.agensi.io/learn/codex-cli-agents-md-complete-guide}{\beamergotobutton{Codex CLI guide}}
\end{frame}

\begin{frame}{Skills: what they are}
  \begin{itemize}
    \item A folder with a \texttt{SKILL.md}.
    \item Like an onboarding guide for a repeatable task
          (e.g.\ ``format regression tables'').
    \item Metadata always loaded ($\sim$100 tokens);
          instructions load on each prompt; rest only when used.
    \item The \texttt{description} decides when Claude/Codex uses the skill.
  \end{itemize}
  \vfill
  \href{https://platform.claude.com/docs/en/agents-and-tools/agent-skills/overview}{\beamergotobutton{Agent Skills docs}}
\end{frame}

\begin{frame}{Where to find \& install them}
  \begin{itemize}
    \item First-party (trusted): Anthropic's repo + pre-built doc skills
          (pptx, xlsx, docx, pdf).
    \item Community directories (vet first): awesome-claude-skills lists,
          marketplace sites.
    \item Install by surface: Claude Code = drop in
          \texttt{\textasciitilde/.claude/skills/} or \texttt{.claude/skills/};
          claude.ai = upload zip; API = \texttt{/v1/skills}.
    \item Don't sync across surfaces; API sandbox has no network access.
  \end{itemize}
  \vfill
  \href{https://github.com/anthropics/skills}{\beamergotobutton{anthropics/skills}}\quad
  \href{https://github.com/ComposioHQ/awesome-claude-skills}{\beamergotobutton{awesome-claude-skills}}
\end{frame}

\begin{frame}{Are they safe?}
  \begin{itemize}
    \item Anthropic: use skills only from trusted sources --- treat like installing software.
    \item Research: prompt injection in $\sim$36\% of skills studied;
          $\sim$91\% of malicious ones bundle malware too.
    \item Audit every file (incl.\ scripts + secondary refs); flag external URL
          fetches; scope Claude's permissions.
    \item Not covered by zero-data-retention --- relevant for confidential microdata.
  \end{itemize}
  \vfill
  \href{https://snyk.io/blog/toxicskills-malicious-ai-agent-skills-clawhub/}{\beamergotobutton{Snyk ToxicSkills}}\quad
  \href{https://securitylabs.datadoghq.com/articles/malicious-skills-supply-chain-risks-in-coding-agents-with-dynamic-context/}{\beamergotobutton{Datadog Security Labs}}
\end{frame}

\begin{frame}{MCPs in empirical research (I)}
  \textbf{What they are, briefly:} standardized connectors that let Claude
  reach your tools and data (databases, files, APIs, reference managers).
  \begin{itemize}
    \item \textbf{Data access:} pull series directly --- e.g.\ FRED MCP exposes
          800{,}000+ economic time series to Claude.
    \item \textbf{Literature:} search arXiv / OpenAlex / Semantic Scholar and
          push hits into Zotero for review.
    \item \textbf{Your project:} filesystem / database servers so Claude reads
          your data schema and files in place.
  \end{itemize}
  \vfill
  \href{https://github.com/stefanoamorelli/fred-mcp-server}{\beamergotobutton{FRED MCP}}\quad
  \href{https://github.com/54yyyu/zotero-mcp}{\beamergotobutton{Zotero MCP}}\quad
  \href{https://registry.modelcontextprotocol.io/}{\beamergotobutton{MCP Registry}}
\end{frame}

\begin{frame}{MCPs in empirical research (II)}
  \textbf{Workflow \& caution:}
  \begin{itemize}
    \item Chain steps: fetch data $\to$ clean $\to$ estimate $\to$ draft ---
          without leaving the agent.
    \item Reproducibility: version control + issue trackers via MCP keep an
          audit trail of what changed.
    \item \textbf{Security:} a server is code with data + execution access;
          prefer local servers you control, scope credentials, vet third-party ones.
    \item For confidential microdata: avoid remote servers; treat community
          servers like any unvetted dependency.
  \end{itemize}
  \vfill
  \href{https://owasp.org/www-community/attacks/MCP_Tool_Poisoning}{\beamergotobutton{Tool poisoning (OWASP)}}\quad
  \href{https://github.com/hanlulong/awesome-ai-for-economists}{\beamergotobutton{AI for economists}}
\end{frame}

\end{document}
